%% ma: L10-B12-0070
%% lop: 10
%% chuong: 5
%% bai: 12
%% dang: 10.12.7
%% loai: TLN
%% muc: VD
%% dap_an: "0,8"
%% nguon: "chuyen-de-cac-so-dac-trung-cua-mau-so-lieu-khong-ghep-nhom-toan-10-chuong-trinh-moi.pdf (GV. Phan Nhật Linh) (P01), Bài 01 Phần III Câu 4"
%% kien_thuc: "Lớp 10 Bài 12 (số gần đúng, sai số tuyệt đối, độ chính xác, sai số tương đối, quy tròn số)"
%% trang_thai: chinh-thuc
%% ghi_chu: "Tự duyệt theo yêu cầu giáo viên 22/10/2026."
\begin{ex}
Bạn Châu đo hai cạnh của một tấm nhôm hình chữ nhật được $17\pm0{,}01$ mm và $23\pm0{,}01$ mm. Giá trị đúng của diện tích tấm nhôm thuộc một đoạn có độ dài bằng bao nhiêu mm$^2$? (Kết quả làm tròn đến hàng phần mười.)
\shortans{0,8}
\loigiai{Viết hai cạnh là $17+d_1$, $23+d_2$ với $|d_1|,|d_2|\le0{,}01$. Khi đó $S=391+17d_2+23d_1+d_1d_2$ và $|17d_2+23d_1+d_1d_2|\le0{,}4001$.

Diện tích đúng thuộc đoạn $[391-0{,}4001;391+0{,}4001]$, có độ dài $0{,}8002\approx0{,}8$.}
\end{ex}
